\section{Introduction}

The FFN sublayer dominates a large fraction of decoder-only transformer
arithmetic. Contextual sparsity observes that the useful FFN components vary
with the input, suggesting that inference can skip work. A FLOP count alone is
not a CPU performance model: selecting individual neurons produces indirect
accesses, fragmented weight reads, branches, and small operations that compete
with cache lines, prefetchers, and vector units.

ARM-Sparse studies a specific compromise. Intermediate neurons are partitioned
into fixed contiguous blocks. A selected block packs corresponding gate and up
rows with down-projection columns and is reused without per-token repacking.
The executor may calculate unwanted neurons inside a block, but it obtains
regular addressing and simpler dispatch. The central hypothesis is that some
block size can outperform theoretically sparser neuron-level execution.

This report does not present contextual sparsity or neuron grouping as new.
Its contribution is an evidence-preserving investigation of the joint
quality--latency break-even on a constrained ARM laptop. The result is
negative but informative: both tested model geometries reach an executor
speedup only well beyond their measured quality-compatible sparsity.
